% ==========================================================
%  Problem Set 5  –  Exercise 1  •  Beamer presentation
%  E383 – Monetary Policy Shocks in the UK
%
%  This template fulfils the tasks in PS5:
%    1. Title page                      (slide: titlepage)
%    2. Outline                         (slide: outline)
%    3. Identification strategy         (formula + data + shocks plot)
%    4. Dynamic causal effects          (LP formula + IRFs)
%    5. Conclusion                      (slide: conclusion)
%  Extra: bibliography frame.
% ==========================================================

\documentclass[aspectratio=169]{beamer}          % widescreen 16:9
% \documentclass[notes=show]{beamer}              % uncomment to show notes

% ---------- Packages ----------
\usetheme[progressbar=frametitle,block=fill]{metropolis}   % modern theme
\usepackage{appendixnumberbeamer}              % number frames correctly
\usepackage{graphicx, subcaption}              % figures & subfigures
\usepackage{svg, tikz}                         % vector graphics
\usetikzlibrary{positioning}
\usepackage{booktabs, threeparttable, multicol} % tables
\usepackage{minted}                            % code listings (needs -shell-escape)
\usepackage{hyperref, marvosym}                % hyperlinks, symbols
\usepackage[english]{babel}                    % language/ hyphenation

% % Uncomment one of the commands to use the University of Tübingen colorcode specified in the style (.sty) file
% \usecolortheme{Tuebingen}
% \input{beamercolorthemeTuebingen.sty}

% ---------- Bibliography ----------
\usepackage[
  backend=biber,
  style=authoryear-comp,
  doi=false,isbn=false,url=false,eprint=false,
  maxcitenames=2,uniquename=false
]{biblatex}
\addbibresource{bibliography.bib}
\setbeamertemplate{bibliography item}{}        % small indent

% ---------- Meta data ----------
\title[PS5]{Problem Set 5}
\subtitle{Effects of Monetary Policy in the UK}
\author[Yannik Winkelmann]{Yannik Winkelmann \\ University of Tübingen}
\date{E383 – SoSe 2025}

% ==========================================================
\begin{document}

% ---------- Title slide ----------
\maketitle

% ---------- Outline ----------
\begin{frame}{Outline}
  \tableofcontents  % auto-populated from \section commands below
\end{frame}

% ==========================================================
\section{Identification Strategy}

% ---------- Identification procedure summary ----------
\begin{frame}{Identification strategy by \textcite{cloyneMacroeconomicEffectsMonetary2016}}
\textcite{cloyneMacroeconomicEffectsMonetary2016} need to tackle the following issues in the identification of monetary policy shocks:
\begin{enumerate}
    \item Simultaneity of monetary policy instruments, interest rates and other macroeconomic variables.
    \item The Bank of England (BoE) builds their decision on a large information set that contains past and expected future economic conditions.
    \item The decisions are based on real-time data and not ex-post data.
\end{enumerate}
\end{frame}

\begin{frame}{Identification strategy by \textcite{cloyneMacroeconomicEffectsMonetary2016}}
How they tackle these issues:
\begin{enumerate}
    \item They utilize the BoE's policy rate - the Bank Rate - as the indented policy target (In contrast \textcite{romer2004new} need to construct the target federal funds rate from minutes of the FOMC meetings). 
    \item They proxy the BoE's information set by constructing an extensive new dataset of historical Bank of England forecasts, private sector forecasts, and real-time data.
    \item The purge this target series of discretionary policy changes that were responding to the changes in macroeconomic variables within the policymakers’ information set.
    
\end{enumerate}
\end{frame}


% ---------- Regression formula ----------
\begin{frame}{Identification strategy – regression}
\textcite{cloyneMacroeconomicEffectsMonetary2016} identify UK monetary-policy shocks as the residual $\varepsilon_m$ from the following regression:
\begin{align}
  \Delta i_m ={}& \alpha  
  + \beta i_{m-d14}
  + \sum_{j=-1}^{2} \gamma_j \hat{y}_{m,j}^{F}
  + \sum_{j=-1}^{2} \varphi_j \hat{\pi}_{m,j}^{F} \notag\\
  &+ \delta \bigl(\hat{y}_{m,j}^{F}-\hat{y}_{m-1,j}^{F}\bigr)
  + \delta \bigl(\hat{\pi}_{m,j}^{F}-\hat{\pi}_{m-1,j}^{F}\bigr)
  + \sum_{j=1}^{3} \rho_j u_{m-j} + \varepsilon_m,  \label{eq:identification}
\end{align}
where the subscript $m$ descries the measurement at meeting-by-meeting frequency, the subscript $j$ denotes the quarter of the forecast relative to the meeting date and the superscript $F$ indicates that the variables are forecasts.
\end{frame}

% ---------- Regression results ----------
\begin{frame}{LaTeX tables with statsmodels}
    \inputminted{python}{latex_table.py}
\end{frame}

\begin{frame}{Regression results}

        \begin{table}[]
            \begin{center}
            \caption{Dependent variable: Change in policy target around the policy decision (FDBRD14)}
            \begin{tabular}{lcccc}
        \toprule \hline
        & \textbf{Estimate} & \textbf{std err} & \textbf{t} & \textbf{P$> |$t$|$}  \\
        \midrule
        \textbf{Intercept}  &      -0.1772  &        0.279     &    -0.634  &         0.526         \\
        $i_{m-d14}$    &      -0.0021  &        0.026     &    -0.081  &         0.935        \\
        $\hat{y}_{m,0}^F$    &       0.0726  &        0.041     &     1.767  &         0.079        \\
        $\hat{y}_{m,1}^F$    &       0.0495  &        0.047     &     1.043  &         0.298      \\
        ... & & & &  \\
        % \bottomrule
        \end{tabular}
            \begin{tabular}{lclc}
            \toprule
            \textbf{No. Observations:} & 235 & \textbf{  R-squared:         } & 0.291 \\
            \textbf{  F-statistic:       } & 4.399 & \textbf{  Prob (F-statistic):} &  1.42e-08  \\
            \bottomrule
            \end{tabular}
            \end{center}

        \end{table}    
\end{frame}

\begin{frame}{Shocks}
    \begin{figure}
        \centering
        \includegraphics[scale=0.5]{figs/shocks.png}
        \caption{Series of monthly monetary policy shocks in the UK from 1975:03 - 2007:05. The shocks are given as the residuals from estimating (\ref{eq:identification}).}
        \label{fig:shocks}
    \end{figure}
\end{frame}

% ==========================================================
\section{Dynamic Causal Effects}

% ---------- Local Projection formula ----------
\begin{frame}{Local-projection regression}
Following \textcite{jorda2005estimation}, the dynamic effect of a one-off shock is estimated using the following LP
\begin{equation}
  y_{t+h}-y_t \;=\; \alpha_h + \theta_h \,\textit{shock}_t 
                   + \beta_h^{\prime}\mathbf X_{t-1} + \varepsilon_{t+h},
  \label{eq:LP}
\end{equation}
where $\text{shock}$ is the monetary policy shock from ($\ref{eq:identification}$), $\mathbf X_{t-1}$ stacks \emph{four lags} of industrial production, CPI, commodity price index and 48 lags of the shocks.
\end{frame}

% ---------- IRF plots ----------
\begin{frame}{Impulse Response Functions}
    \begin{figure}
        \centering
        \includegraphics[scale=0.225]{figs/IRFs.png}
        \caption{\small Impuse respones given a 1\% contractionary monetary policy shock as estimated by (\ref{eq:LP}). Confidence bands indicate the 90 percent interval with standard errors are adjusted for serial correlation. \normalsize}
        \label{fig:IRFsl}
    \end{figure}
\end{frame}


% ==========================================================
\section{Conclusion}

\begin{frame}{Conclusion}
\begin{itemize}
  \item Narrative identification à la \textcite{cloyneMacroeconomicEffectsMonetary2016}
        isolates unexpected changes in the UK policy rate.
  \item Local projections reveal:
        \begin{itemize}
            \item Economy slows down following the unexpected increase in the policy rate.
            \item CPI declines significantly only after $\sim$30 months.
        \end{itemize}
\end{itemize}
\end{frame}

% ==========================================================
\begin{frame}[allowframebreaks]{References}
  \printbibliography[heading=none]
\end{frame}

\end{document}
